\documentclass[12pt,a4paper]{article}
\usepackage[utf8]{inputenc}
\usepackage[T1]{fontenc}
\usepackage{amsmath,amssymb,graphicx,booktabs,hyperref,natbib}
\usepackage[margin=2.5cm]{geometry}

\begin{document}

\title{\bf A non-monotonic dark energy equation of state\\
derived from the cosmic star formation history}

\author{Anonymous}
\date{June 2026}

\maketitle

\begin{abstract}
The physical nature of dark energy is unknown. The DESI collaboration has reported evidence that the dark energy equation of state evolves with time, but all existing models treat $w(z)$ as a function to be fitted rather than derived from independent astrophysical measurements. Here we show that if cosmic structure formation occupies a finite reservoir of vacuum quantum modes, the dark energy equation of state acquires a specific, non-monotonic shape determined by the cosmic star formation history: $w(z)+1\propto{\rm SFR}(z)/[H(z)\,\rho_*(z)^{n/(n+1)}]$. We solve the model self-consistently by iterating the coupled Friedmann and mode-occupation equations, finding that negative feedback stabilizes the solution near $\Lambda$CDM with $H(z)$ deviating by approximately $2\%$. The converged prediction is a non-monotonic $w(z)$: $w_0=-0.80$, rising to a peak of $w\approx-0.31$ at $z\approx1.5$, then returning toward $-1$ at higher redshift. Against DESI DR2 baryon acoustic oscillation distances, the model yields $\chi^2=19.1$ for 12 data points with two parameters, compared to $\chi^2=15.0$ for $\Lambda$CDM. The model is not statistically preferred over a cosmological constant on current data. Its distinguishing feature---a peaked equation of state---cannot be produced by the CPL parameterization or by quintessence with simple potentials for any choice of their parameters. We show that the SFR input is robust against the change in luminosity distances induced by the model, and that DESI DR3 and Euclid will distinguish the predicted peak from a monotonic function at $3.4\sigma$. If $w(z)$ is monotonic, the model is falsified.
\end{abstract}

\section*{Introduction}

A quarter-century after the discovery of cosmic acceleration, the physical nature of dark energy remains unknown. The cosmological constant $\Lambda$ is consistent with most cosmological data but offers no mechanism, no dynamics, and no connection to the rest of physics. Dynamical dark energy models---quintessence~\cite{tsujikawa2013}, emergent dark energy~\cite{li2024}, and phenomenological parameterizations~\cite{chevallier2001}---allow the equation of state $w(z)$ to vary but share a common structure: they choose a functional form for $w(z)$ by hand, fit its parameters to data, and provide no answer to the question of why that particular form should describe nature.

The DESI collaboration has recently sharpened this question. Using baryon acoustic oscillations measured from over 14 million galaxies, DESI DR2 finds a $2.8$--$4.2\sigma$ preference for dynamical dark energy over $\Lambda$CDM~\cite{desi2025}. In the CPL parameterization, the best fit is $(w_0,w_a)=(-0.785\pm0.047,-0.43\pm0.10)$ for DESI+CMB+DESY5~\cite{li2026}. But CPL describes {\it that} $w$ evolves, not {\it why}.

Here we derive $w(z)$ from a specific physical hypothesis: cosmic structure formation occupies a finite reservoir of vacuum quantum modes, and the dark energy density is the energy of the modes that remain free. The idea that spacetime possesses a finite information capacity has a long history, from the holographic principle~\cite{thooft1993,susskind1995} through causal set theory~\cite{sorkin2005} to entropic gravity~\cite{verlinde2011,jacobson1995}. In the framework we use, the fraction $q(t)\in[0,1]$ of unoccupied modes obeys a telegraph equation sourced by the cosmic star formation rate. The dark energy density is $\rho_{\rm DE}(t)=\rho_\Lambda q(t)$, and $w(z)$ is determined by the star formation history---an independently measured quantity.

This is a phenomenological model: the coupling between star formation and mode occupation is calibrated from $w_0$, not derived from first principles. What makes it different from existing approaches is the structure of its prediction. Because $w(z)+1$ is proportional to the star formation rate divided by a monotonic function of the cumulative stellar mass, it inherits the peak that the cosmic star formation history exhibits at $z\sim1$--$2$~\cite{madau2014}. This non-monotonic shape cannot be produced by CPL or by quintessence with simple potentials. It is a falsifiable consequence of the hypothesis.

\section*{Model and self-consistent solution}

The occupied fraction $\Delta\equiv1-q$ obeys a nonlinear telegraph equation. On homogeneous cosmological scales, with the memory term set to zero (the data prefer this limit; Supplementary Information):

\begin{equation}\label{eq:ode}
\gamma_0\Delta^n\frac{d\Delta}{dt} = \eta\,\dot{\rho}_*(t),
\end{equation}

where $\dot{\rho}_*(t)={\rm SFR}(t)$ is the cosmic star formation rate density, $\gamma_0$ is a damping rate, $n$ is a damping nonlinearity index, and $\eta$ is a coupling constant. Integration yields $\Delta^{n+1}\propto\rho_*(t)$, with $\rho_*$ the cumulative stellar mass density. The dark energy density is $\rho_{\rm DE}=\rho_\Lambda(1-\Delta)$. From energy conservation,

\begin{equation}\label{eq:shape}
w(z)+1 = C\cdot\frac{{\rm SFR}(z)}{H(z)\,\rho_*(z)^{\,n/(n+1)}}\cdot\frac{1}{1-\Delta(z)},
\end{equation}

where $C\propto\sqrt{\eta/\gamma_0}$ is calibrated from $w_0$. The natural value $n=1$ corresponds to damping proportional to the information deficit. A full scan over $n\in[0.3,3.0]$ and $w_0\in[-0.95,-0.65]$ changes $\chi^2$ by $<3$ units (Supplementary Information); the prediction is robust.

Equation (\ref{eq:shape}) depends on $H(z)$ on both sides because $\rho_*(z)=\int_z^\infty{\rm SFR}(z')[(1+z')H(z')]^{-1}dz'$. We solve by Picard iteration: starting from $\Lambda$CDM $H(z)$, compute $\rho_*$, $w$, and $\Delta$; update $H(z)$ via the Friedmann equation; recompute $\rho_*$; repeat. The iteration converges in $<15$ steps to $\max|\delta H/H|<10^{-5}$, driven by negative feedback: a larger $H$ reduces $|dt/dz|$, reducing $\rho_*$, bringing $H$ back toward $\Lambda$CDM. This feedback is a physical feature of the coupled system, not a numerical artifact: structure formation affects expansion, and expansion affects the rate at which structure forms.

We use the Madau \& Dickinson (2014) SFR fit~\cite{madau2014} and Planck 2018 parameters~\cite{planck2018}, with the sound horizon $r_d$ as a free parameter. The SFR fit was calibrated from UV and IR photometry assuming $\Lambda$CDM luminosity distances. Because DGF modifies $H(z)$ by only $\sim2\%$, the resulting change in luminosity distances produces a $3$--$5\%$ correction to the inferred SFR, which propagates to a $<0.5\%$ change in $w(z)$---well below the $\pm27\%$ systematic uncertainty in the SFR normalization (Supplementary Information). The model is therefore robust against the cosmological circularity of its SFR input.

\section*{Comparison with DESI DR2}

We compare the converged prediction to DESI DR2 BAO distances~\cite{desi2025} at six effective redshifts (12 data points: $D_M/r_d$ and $D_H/r_d$ each), including the anti-correlation $\rho\approx-0.5$ between the two distance measures.

\begin{table}[htbp]
\centering
\caption{Model comparison against DESI DR2 BAO distances ($r_d$ free).}
\label{tab:models}
\begin{tabular}{lccccc}
\toprule
Model & DE params & $r_d$ (Mpc) & $\chi^2$ & AIC & \\
\midrule
$\Lambda$CDM     & 0 ($w\equiv-1$)   & 148.6 & 15.0 & 17.0 \\
DGF $n=1$       & 2 ($C$, $n$)      & 146.0 & 19.1 & 25.1 \\
DGF $n=2$       & 2                  & 146.7 & 17.7 & 23.7 \\
\bottomrule
\end{tabular}
\end{table}

Table~\ref{tab:models} reports the results. With $r_d$ free, $\Lambda$CDM achieves $\chi^2=15.0$ (best-fit $r_d=148.6$~Mpc). DGF yields $\chi^2=19.1$ ($n=1$) or $17.7$ ($n=2$), with AIC values of 25.1 and 23.7, compared to 17.0 for $\Lambda$CDM. The model is not statistically preferred over the cosmological constant. A full scan over the parameter space $(w_0,n,r_d)$ yields no combination with AIC below that of $\Lambda$CDM; the $\sim2\%$ deviation of $H(z)$ from $\Lambda$CDM, stabilized by the negative feedback mechanism, is a structural feature of the coupled system that DESI BAO precision is sufficient to penalize.

We do not regard this as a failure of the model. The purpose of DGF is not to provide a superior fit to BAO data---a task for which the two-parameter CPL form is already adequate. Its purpose is to make a specific, falsifiable prediction that no other model makes. The $\chi^2$ comparison establishes that the model is not excluded by current data. The prediction is what distinguishes it.

\section*{A non-monotonic prediction}

Figure~\ref{fig:wz} shows the converged $w(z)$. The key feature is the peak: $w$ rises from $-0.80$ at $z=0$ to approximately $-0.31$ at $z\approx1.5$, then falls back toward $-1$ at higher redshift. This shape is inherited from the cosmic star formation history, which rises from the present epoch to a maximum at $z\approx1.9$~\cite{madau2014} before declining. The peak amplitude is moderated by negative feedback in the self-consistent iteration: without it, the driving-dominated approximation would produce a much larger deviation (Supplementary Information).

\begin{figure}[htbp]
\centering
\includegraphics[width=0.85\textwidth]{../../LP32-S8_DESI-Contour/figures/fig2_wz_shape.png}
\caption{Self-consistently converged $w(z)$ for $n=1$ (blue), with $n=0.7$ and $n=1.3$. The peak is the model's distinguishing prediction. The band shows the propagated $\pm27\%$ SFR normalization uncertainty.}
\label{fig:wz}
\end{figure}

Table~\ref{tab:competitors} compares this prediction to existing dark energy models. The CPL parameterization $w(a)=w_0+w_a(1-a)$ is linear in $(1-a)$ and monotonic. Quintessence models with power-law potentials $V(\phi)\propto\phi^{-\alpha}$ are also monotonic~\cite{tsujikawa2013}. The emergent dark energy model~\cite{li2024} can produce non-monotonic behavior in some parameter regimes but lacks a connection to the star formation history; its $w(z)$ shape is chosen, not derived from independent data. No existing model with a comparable number of parameters makes the specific prediction that $w(z)$ inherits the shape of the cosmic star formation history.

\begin{table}[htbp]
\centering
\caption{Comparison of dark energy models.}
\label{tab:competitors}
\begin{tabular}{lccccc}
\toprule
Model & $N_{\rm par}$ & $w(z)$ source & Non-monotonic? & SFR connection? \\
\midrule
$\Lambda$CDM & 0 & Fixed ($w\equiv-1$) & No & No \\
CPL & 2 & Phenomenological & No & No \\
Quintessence & 2--3 & Scalar field potential & No & No \\
Emergent DE & 2--3 & Integration constant & Possible & No \\
{\bf This work} & {\bf 2} & {\bf Derived from SFR} & {\bf Yes} & {\bf Yes} \\
\bottomrule
\end{tabular}
\end{table}

\section*{Testability}

The peak in $w(z)$ is a falsifiable prediction. Using a Fisher matrix forecast for five redshift bins of width $\Delta z=0.5$, including cross-bin correlations of $\rho=0.4$ between adjacent bins, we find that DESI DR3 and Euclid will distinguish the DGF peak from the CPL best-fit monotonic function at $3.4\sigma$ (combined). DESI DR3 alone will provide approximately $2.3\sigma$. The sensitivity is concentrated at $z\sim1$--$2$, where the DGF peak is most pronounced and where the CPL and quintessence predictions are monotonic. If $w(z)$ is found to be consistent with a linear function in $(1-a)$, the model is excluded.

The model is also consistent with existing cosmic microwave background constraints. Because DGF modifies only the late-time ($z<5$) expansion history, its predictions for the CMB are identical to those of $\Lambda$CDM at recombination. The mild tension between DESI BAO distances and the Planck-calibrated sound horizon---which affects $\Lambda$CDM at the $\sim2\sigma$ level---is not exacerbated by DGF.

\section*{Discussion}

The purpose of a physical theory is not to fit existing data better than a phenomenological parameterization---it is to make predictions that can be wrong. The CPL parameterization, for all its utility, cannot be wrong: any monotonic $w(z)$ can be accommodated by adjusting $(w_0,w_a)$. Quintessence, for all its physical motivation, cannot be wrong in the same sense: a suitably chosen potential can approximate any smooth $w(z)$. The model presented here is different. It predicts a specific, non-monotonic shape for $w(z)$ that follows from the independently measured cosmic star formation history. If future data show $w(z)$ to be monotonic, the hypothesis is falsified.

We have been explicit about the model's limitations. The coupling between star formation and mode occupation is calibrated from $w_0$, not computed from first principles. The damping index $n$ is motivated by physical reasoning but not uniquely derived. The model does not outperform $\Lambda$CDM on current BAO data. These are genuine limitations, and we state them without reservation.

What the model provides is a connection---tentative, falsifiable, and quantitative---between dark energy and the cosmic history of structure formation. It will be tested.

\begin{thebibliography}{99}
\bibitem{riess1998} Riess, A.G.\ et al.\ {\it Astron.\ J.} {\bf 116}, 1009 (1998).
\bibitem{spergel2003} Spergel, D.N.\ et al.\ {\it Astrophys.\ J.\ Suppl.} {\bf 148}, 175 (2003).
\bibitem{tsujikawa2013} Tsujikawa, S.\ {\it Class.\ Quant.\ Grav.} {\bf 30}, 214003 (2013).
\bibitem{li2024} Li, X.\ et al.\ {\it Phys.\ Rev.\ D} {\bf 109}, 043510 (2024).
\bibitem{chevallier2001} Chevallier, M.\ \& Polarski, D.\ {\it Int.\ J.\ Mod.\ Phys.\ D} {\bf 10}, 213 (2001).
\bibitem{desi2025} DESI Collaboration.\ {\it Phys.\ Rev.\ D} {\bf 112}, 083515 (2025).
\bibitem{li2026} Li, T.-N.\ et al.\ {\it Phys.\ Dark Universe} (2026); arXiv:2511.22512.
\bibitem{thooft1993} 't Hooft, G.\ arXiv:gr-qc/9310026 (1993).
\bibitem{susskind1995} Susskind, L.\ {\it J.\ Math.\ Phys.} {\bf 36}, 6377 (1995).
\bibitem{sorkin2005} Sorkin, R.D.\ In {\it Approaches to Quantum Gravity} (Cambridge, 2009).
\bibitem{verlinde2011} Verlinde, E.\ {\it J.\ High Energy Phys.} {\bf 2011}, 29 (2011).
\bibitem{jacobson1995} Jacobson, T.\ {\it Phys.\ Rev.\ Lett.} {\bf 75}, 1260 (1995).
\bibitem{madau2014} Madau, P.\ \& Dickinson, M.\ {\it Annu.\ Rev.\ Astron.\ Astrophys.} {\bf 52}, 415 (2014).
\bibitem{planck2018} Planck Collaboration.\ {\it Astron.\ Astrophys.} {\bf 641}, A6 (2020).
\end{thebibliography}

\end{document}
